\documentclass[11pt]{article}
\usepackage{amsmath,amssymb,amsthm}
\usepackage{geometry}
\geometry{margin=1in}
\usepackage{hyperref}
\usepackage{booktabs}
\usepackage{microtype}

% Theorem environments
\newtheorem{theorem}{Theorem}[section]
\newtheorem{proposition}[theorem]{Proposition}
\newtheorem{corollary}[theorem]{Corollary}
\newtheorem{lemma}[theorem]{Lemma}
\theoremstyle{definition}
\newtheorem{definition}[theorem]{Definition}
\theoremstyle{remark}
\newtheorem{remark}[theorem]{Remark}

% Convenience macros
\newcommand{\Zfi}{\mathbb{Z}[\varphi]}
\newcommand{\Qfi}{\mathbb{Q}(\sqrt{5})}
\newcommand{\hqa}{H_{\mathrm{QA}}}
\newcommand{\hraw}{H_{\mathrm{raw}}}
\newcommand{\Nm}[1]{N\!\left(#1\right)}

\title{\textbf{The Arithmetic of the QA Map:\\
$\mathbb{Z}[\varphi]$, Norm Invariants, and Half-Pisano Orbits}}
\author{[Author redacted for review]}
\date{}

\begin{document}
\maketitle

\begin{abstract}
The Quantum Arithmetic (QA) map $T:(b,e)\mapsto(b{+}e,\,b{+}2e)\pmod{m}$
is the linear representation of multiplication by $\varphi^2$ in the ring
$\Zfi$ of integers of $\Qfi$, where $\varphi=(\sqrt{5}+1)/2$.
The quadratic form $f(b,e)=b^2+be-e^2$, which is the algebraic norm
$N(b+e\varphi)$ in $\Qfi$, is $T$-invariant; mod-$m$ values of $f$ classify
QA orbits. Since $T=Q^2$ where $Q$ is the Fibonacci companion matrix, the
orbit length of any primitive state equals $\pi(m)/2$, half the Pisano period
(Wall, 1960). The projective ratio $z_t=e_t/(b_t+e_t)$ evolves by a
M\"{o}bius map with unique attracting fixed point $z^*=1/\varphi$.
Finally, the harmonic index $\hqa$ used in QA-based learning admits a clean
operator decomposition: $4\hraw=\mathrm{tr}(\Lambda)$ where
$\Lambda=\mathrm{diag}(\sigma_d,\sigma_{od})$ is a $2\times 2$ cross-coupling
operator; the equal-weight Rayleigh identity $2\hraw=R(\Lambda,(1,1)^\top/\sqrt{2})$
is canonical by a three-part symmetry argument.
This note consolidates the algebraic foundations; optimization applications
appear in the companion paper~\cite{companion}.
\end{abstract}

\tableofcontents
\bigskip

%% ────────────────────────────────────────────────────────────────────────────
\section{Introduction}
%% ────────────────────────────────────────────────────────────────────────────

The Quantum Arithmetic (QA) system is a modular arithmetic framework whose
dynamics arise from a simple recurrence: given a pair $(b,e)\in\{1,\ldots,m\}^2$,
one step produces $(d,a)=(b+e\bmod m, b+2e\bmod m)$, with the convention that
$0$ is replaced by $m$.  The resulting sequences are Fibonacci-like modulo $m$,
and the QA map $T:(b,e)\mapsto(d,a)$ generates a partition of the state space
into finitely many closed orbits.

At first glance this is a piece of recreational combinatorics.  The purpose of
this note is to show that the structure is, in fact, determined entirely by the
arithmetic of $\Qfi$: the map is multiplication by $\varphi^2$, the invariant
of each orbit is the algebraic norm, and orbit lengths follow from the Pisano
period.  These are classical facts in algebraic number theory, here made explicit
for the QA setting.

The note is organised as follows.  Section~\ref{sec:ring} defines the ring
$\Zfi$ and identifies $T$ as $Q^2$.  Section~\ref{sec:norm} proves norm
invariance and gives the geometric determinant interpretation.
Section~\ref{sec:pisano} recalls the Half-Pisano theorem and derives orbit
lengths.  Section~\ref{sec:projective} analyses the projective Fibonacci flow.
Section~\ref{sec:hqa} derives the $\hqa$ operator decomposition.
Section~\ref{sec:conclusion} concludes with directions for the companion paper.

%% ────────────────────────────────────────────────────────────────────────────
\section{The Ring $\Zfi$ and the QA Map}
\label{sec:ring}
%% ────────────────────────────────────────────────────────────────────────────

\subsection{The golden ratio and its ring of integers}

The golden ratio $\varphi=(\sqrt{5}+1)/2$ satisfies $\varphi^2=\varphi+1$ and is
a root of the minimal polynomial $x^2-x-1$.  Its Galois conjugate is
$\bar\varphi=(1-\sqrt{5})/2\approx-0.618$, satisfying $\varphi+\bar\varphi=1$ and
$\varphi\bar\varphi=-1$.

The ring of integers of $\Qfi$ is $\Zfi=\{a+b\varphi:a,b\in\mathbb{Z}\}$, with
addition and multiplication inherited from $\mathbb{R}$.  Multiplication by
$\varphi$ acts on the coordinate pair $(a,b)$ as:
\[
  \varphi(a+b\varphi)=a\varphi+b\varphi^2=a\varphi+b(\varphi+1)=b+(a+b)\varphi,
\]
so the matrix of ${\times}\varphi$ in the basis $\{1,\varphi\}$ is
\[
  Q = \begin{pmatrix}0&1\\1&1\end{pmatrix}.
\]
This is the Fibonacci companion matrix: if $(F_n)$ denotes the Fibonacci
sequence then $Q^n=\bigl(\begin{smallmatrix}F_{n-1}&F_n\\F_n&F_{n+1}
\end{smallmatrix}\bigr)$.

\subsection{The QA map as $Q^2$}

\begin{proposition}[QA map is $\times\varphi^2$]
\label{prop:qa-is-phi2}
Let $T:(b,e)\mapsto(b{+}e,\,b{+}2e)\pmod{m}$ be the QA map (with $0$ replaced
by $m$).  Then $T=Q^2$ as a map on $(\mathbb{Z}/m\mathbb{Z})^2$, and corresponds
to multiplication by $\varphi^2=\varphi+1$ in $\Zfi/m\Zfi$.
\end{proposition}

\begin{proof}
Direct computation:
\[
  Q^2 = \begin{pmatrix}0&1\\1&1\end{pmatrix}^2
      = \begin{pmatrix}1&1\\1&2\end{pmatrix}.
\]
Applied to $(b,e)^\top$: $Q^2(b,e)^\top=(b+e,\,b+2e)^\top=(d,a)^\top$.
The ring interpretation: $(b+e\varphi)\cdot\varphi^2=(b+e\varphi)(\varphi+1)$
$=b\varphi+b+e\varphi^2+e\varphi=b+e+(b+2e)\varphi=d+a\varphi$.
\end{proof}

\begin{remark}
The map $Q^1$ would send $(b,e)\mapsto(e,b+e)$, i.e., one Fibonacci step.
$T=Q^2$ skips the intermediate pair and maps directly to the pair two steps ahead.
The QA system thus ``runs the Fibonacci recurrence at double speed''.
\end{remark}

%% ────────────────────────────────────────────────────────────────────────────
\section{The Quadratic Norm and Orbit Invariant}
\label{sec:norm}
%% ────────────────────────────────────────────────────────────────────────────

\subsection{Algebraic norm}

The algebraic norm in $\Qfi$ is $N(\alpha)=\alpha\bar\alpha$ for
$\alpha=a+b\varphi$:
\[
  \Nm{a+b\varphi}
  = (a+b\varphi)(a+b\bar\varphi)
  = a^2+ab(\varphi+\bar\varphi)+b^2\varphi\bar\varphi
  = a^2+ab-b^2.
\]
For a QA state $(b,e)$ we define $f(b,e)=\Nm{b+e\varphi}=b^2+be-e^2$.

\begin{proposition}[$f$ is $T$-invariant]
\label{prop:norm-invariant}
For any $(b,e)\in\mathbb{Z}^2$, $f(T(b,e))=f(b,e)$.
\end{proposition}

\begin{proof}
Since $T$ is multiplication by $\varphi^2$ and the norm is multiplicative:
\[
  f(T(b,e)) = \Nm{(b+e\varphi)\varphi^2}
            = \Nm{b+e\varphi}\cdot\Nm{\varphi^2}
            = f(b,e)\cdot N(\varphi)^2
            = f(b,e)\cdot(-1)^2
            = f(b,e). \qedhere
\]
\end{proof}

\subsection{Geometric determinant identity}

The norm has a direct geometric meaning in terms of the QA step matrix.

\begin{proposition}[Determinant identity]
\label{prop:det}
Let $M=\bigl(\begin{smallmatrix}b&e\\d&a\end{smallmatrix}\bigr)$ where $(d,a)=T(b,e)$
in $\mathbb{Z}^2$ (before reduction mod $m$).  Then $\det(M)=f(b,e)$.
\end{proposition}

\begin{proof}
$\det(M)=ba-ed=b(b+2e)-e(b+e)=b^2+2be-eb-e^2=b^2+be-e^2=f(b,e)$.
\end{proof}

\subsection{Orbit classification for $m=9$}

For $m=9$ the state space is $\{1,\ldots,9\}^2$.
The 3-adic valuation $v_3(f(b,e)\bmod 9)$ determines the orbit type:

\begin{table}[h]
\centering
\caption{Orbit classification for $m=9$ via the norm $f(b,e)\bmod 9$.}
\begin{tabular}{lll}
\toprule
$v_3(f(b,e))$ & Orbit length & Orbit type \\
\midrule
$\geq 4$ & $1$  & Singularity (fixed point $(9,9)$) \\
$2$--$3$ & $4$  & Satellite \\
$0$      & $12$ & Cosmos (primitive) \\
\bottomrule
\end{tabular}
\label{tab:orbits}
\end{table}

This classification is verified exhaustively for all 81 states of
$\{1,\ldots,9\}^2$.  Primitive orbits (Cosmos) have $\gcd(f(b,e),9)=1$;
there are 72 such states forming six 12-cycles.

%% ────────────────────────────────────────────────────────────────────────────
\section{Orbit Lengths and the Half-Pisano Theorem}
\label{sec:pisano}
%% ────────────────────────────────────────────────────────────────────────────

\subsection{The Pisano period}

The \emph{Pisano period} $\pi(m)$ is the period of the Fibonacci sequence modulo
$m$: $F_{n+\pi(m)}\equiv F_n\pmod{m}$ for all $n\geq 0$.  Equivalently,
$\pi(m)$ is the order of the matrix $Q$ in $\mathrm{GL}_2(\mathbb{Z}/m\mathbb{Z})$.

\begin{theorem}[Half-Pisano Orbit Length, Wall~\cite{wall1960}]
\label{thm:half-pisano}
For $m\geq 3$, $\pi(m)\equiv 0\pmod{2}$.  The order of $T=Q^2$ in
$\mathrm{GL}_2(\mathbb{Z}/m\mathbb{Z})$ is $\pi(m)/2$.  Consequently,
the orbit length of any primitive state $(b,e)$ under the QA map $T$ divides
$\pi(m)/2$, and equals $\pi(m)/2$ if and only if $(b,e)$ is primitive.
\end{theorem}

\begin{proof}
\emph{Step 1 (Pisano period is even).}
Wall~\cite{wall1960} proves that for $m\geq 3$, the Fibonacci sequence
mod $m$ satisfies $F_{\pi(m)-1}\equiv 1\pmod{m}$ and
$F_{\pi(m)}\equiv 0\pmod{m}$, which forces $Q^{\pi(m)}=I$ and, by
analysis of the entry parities, $\pi(m)\equiv 0\pmod 2$.

\emph{Step 2 (Order of $Q^2$).}
$\mathrm{ord}(Q^2)=\mathrm{ord}(Q)/\gcd(2,\mathrm{ord}(Q))=\pi(m)/2$
(since $\pi(m)$ is even).

\emph{Step 3 (Orbit length).}
The orbit of $(b,e)$ under $T$ has length $\mathrm{ord}(Q^2)/\mathrm{stab}$,
where $\mathrm{stab}$ is the size of the stabiliser.  For primitive $(b,e)$
the stabiliser is trivial, giving length $\pi(m)/2$.
\end{proof}

\begin{corollary}[Pisano values for small $m$]
\label{cor:pisano}
\hfill
\begin{center}
\begin{tabular}{ccc}
\toprule
$m$ & $\pi(m)$ & Primitive orbit length $\pi(m)/2$ \\
\midrule
3  & 8  & 4 \\
5  & 20 & 10 \\
7  & 16 & 8 \\
9  & 24 & 12 \\
24 & 24 & 12 \\
\bottomrule
\end{tabular}
\end{center}
\end{corollary}

%% ────────────────────────────────────────────────────────────────────────────
\section{Projective Flow and the Golden Attractor}
\label{sec:projective}
%% ────────────────────────────────────────────────────────────────────────────

\begin{proposition}[Projective Fibonacci flow]
\label{prop:projective}
Let $z_t=e_t/(b_t+e_t)$.  The projective ratio evolves by the M\"{o}bius map
\[
  z_{t+1} = \frac{1+z_t}{2+z_t}
\]
with unique attracting fixed point $z^*=(\sqrt{5}-1)/2=1/\varphi\approx 0.618$.
The convergence rate near $z^*$ is $|z_{t+1}-z^*|\approx|z_t-z^*|/\varphi^2$.
\end{proposition}

\begin{proof}
From $(d,a)=(b+e,b+2e)$:
\[
  z_{t+1}=\frac{e_{t+1}}{d_{t+1}+e_{t+1}}=\frac{a}{d+a}
  =\frac{b+2e}{(b+e)+(b+2e)}=\frac{b+2e}{2b+3e}.
\]
Dividing numerator and denominator by $b+e$ and setting $z=e/(b+e)$:
\[
  z_{t+1}=\frac{1+z}{2+z}.
\]
Fixed points: $z^*=(1+z^*)/(2+z^*)\Rightarrow (z^*)^2+z^*-1=0\Rightarrow
z^*=(\sqrt{5}-1)/2=1/\varphi$.  Linearising about $z^*$: derivative
$(2+z^*-1-z^*)/(2+z^*)^2=1/(2+z^*)^2=1/\varphi^4$...

More directly: near $z^*$, $z_{t+1}-z^*\approx f'(z^*)(z_t-z^*)$ where
$f(z)=(1+z)/(2+z)$ gives $f'(z)=1/(2+z)^2$.  At $z^*=1/\varphi$,
$2+z^*=2+1/\varphi=\varphi+1+1/\varphi=({\varphi^2+1})/\varphi$, and
$f'(z^*)=\varphi^2/(\varphi^2+1)^2$.
Since $\varphi^2=\varphi+1$ and $\varphi^2+1=\varphi+2$,
$f'(z^*)=(\varphi+1)/(\varphi+2)^2<1$, confirming $z^*$ is attracting.
\end{proof}

\begin{remark}
In the unmodded (linear) case, $\sigma_d=\cos(2\arctan(z_t))$ converges to
$\cos(2\arctan(1/\varphi))$ as $t\to\infty$, providing a computable limiting
value for the degree-0 component of $\hraw$.
\end{remark}

%% ────────────────────────────────────────────────────────────────────────────
\section{The $H_{QA}$ Operator Decomposition}
\label{sec:hqa}
%% ────────────────────────────────────────────────────────────────────────────

The harmonic index $\hqa$ used in QA-based optimisation is defined via the
cross-coupling operator of the QA step matrix $M$.

\subsection{The cross-coupling operator}

\begin{definition}[QA step matrix and cross-coupling]
\label{def:hqa}
Given a state $(b,e)$ with one QA step $(d,a)=T(b,e)$ (in $\mathbb{Z}$, before
reduction mod $m$), define the step matrix $M=\bigl(\begin{smallmatrix}b&e\\d&a
\end{smallmatrix}\bigr)$ and the two coupling ratios:
\begin{align*}
  \sigma_d(M)   &= \frac{M_{11}M_{22}}{M_{12}^2+M_{21}^2}
                 = \frac{ba}{e^2+d^2}, \quad\text{(degree-0, scale-invariant)} \\
  \sigma_{od}(M)&= \frac{M_{12}M_{21}}{M_{11}+M_{22}}
                 = \frac{ed}{b+a}. \quad\text{(degree-1)}
\end{align*}
Set $\Lambda=\mathrm{diag}(\sigma_d,\sigma_{od})$ and
$\hraw=\mathrm{tr}(\Lambda)/4$.  The bounded harmonic index is
$\hqa=|\hraw|/(1+|\hraw|)\in(0,1)$.
\end{definition}

\begin{proposition}[Trace and Rayleigh identities]
\label{prop:trace-rayleigh}
\begin{enumerate}[(i)]
  \item $4\hraw = \mathrm{tr}(\Lambda) = \sigma_d+\sigma_{od}$.
  \item $2\hraw = R(\Lambda,v_0)$ where $v_0=(1,1)^\top/\sqrt{2}$
    is the equal-weight unit vector and $R$ denotes the Rayleigh quotient.
  \item $\sigma_{od}=e/2$ via the Fibonacci identity $a+b=2d$.
  \item Both identities satisfy transpose symmetry:
    $\sigma_d(M)=\sigma_d(M^\top)$ and $\sigma_{od}(M)=\sigma_{od}(M^\top)$.
\end{enumerate}
\end{proposition}

\begin{proof}
(i) Immediate from $\mathrm{tr}(\Lambda)=\sigma_d+\sigma_{od}$.
(ii) $R(\Lambda,v_0)=v_0^\top\Lambda v_0=\tfrac{1}{2}(\sigma_d+\sigma_{od})=2\hraw$.
(iii) $a+b=(b+2e)+(b)=2(b+e)=2d$, so $ed/(b+a)=ed/(2d)=e/2$.
(iv) Swapping rows and columns of $M$ sends $(b,e,d,a)\mapsto(b,d,e,a)$, and
$ba/(d^2+e^2)=ba/(e^2+d^2)=\sigma_d$ (unchanged), $de/(b+a)=\sigma_{od}$ (unchanged). \qedhere
\end{proof}

\subsection{Why the equal-weight vector $v_0=(1,1)^\top/\sqrt{2}$}

\begin{proposition}[Canonical Rayleigh vector]
\label{prop:v0}
The vector $v_0=(1,1)^\top/\sqrt{2}$ is the unique unit vector in the positive
orthant satisfying all three of:
\begin{enumerate}[(i)]
  \item \textbf{Swap symmetry.} $v_0$ is invariant under the coordinate permutation
    $(v_1,v_2)\mapsto(v_2,v_1)$, which corresponds to interchanging the two QA
    pairs $(b,e)\leftrightarrow(d,a)$.  Since $T$ is a bijection on every orbit,
    neither pair is algebraically privileged.
  \item \textbf{Transpose consistency.} $\Lambda$ is symmetric under $M\mapsto M^\top$
    (Proposition~\ref{prop:trace-rayleigh}(iv)), so the natural Rayleigh vector
    must also be symmetric.
  \item \textbf{Trace recovery.} $\mathrm{tr}(\Lambda)=2\,R(\Lambda,v_0)$:
    the Rayleigh quotient at $v_0$ recovers the normalised trace of $\Lambda$
    with the unique scalar factor $1/\sqrt{2}$ (up to sign).
\end{enumerate}
\end{proposition}

\begin{proof}
Any unit vector $v=(\cos\theta,\sin\theta)^\top$ with $\theta\in(0,\pi/2)$ satisfies
$R(\Lambda,v)=\sigma_d\cos^2\theta+\sigma_{od}\sin^2\theta$.
Condition (i) requires $\cos\theta=\sin\theta$, hence $\theta=\pi/4$ and
$v=v_0$.  Conditions (ii) and (iii) are then verified directly.
\end{proof}

\subsection{Bilinear identity}

The following polynomial identity relates the two coupling ratios to the orbit
invariant $f(b,e)$ and holds in $\mathbb{Z}$ (before reduction mod $m$):
\begin{equation}
  \label{eq:bilinear}
  4\hraw\cdot(e^2+d^2)\cdot(b+a) = ba(b+a) + ed(e^2+d^2).
\end{equation}

This is an algebraic identity in the four QA values $(b,e,d,a)$, verified
directly by substituting $d=b+e$, $a=b+2e$.  It encodes the balance between
the scale-free ($\sigma_d$) and amplitude ($\sigma_{od}$) components of $\hraw$.

%% ────────────────────────────────────────────────────────────────────────────
\section{Conclusion}
\label{sec:conclusion}
%% ────────────────────────────────────────────────────────────────────────────

We have shown that the QA map $T:(b,e)\mapsto(b+e,b+2e)\pmod{m}$ is, in essence,
the arithmetic of $\Zfi$:

\begin{itemize}
  \item $T=Q^2$ is multiplication by $\varphi^2$ (Proposition~\ref{prop:qa-is-phi2}).
  \item The quadratic norm $f(b,e)=b^2+be-e^2=N(b+e\varphi)$ is $T$-invariant and
    classifies orbits (Propositions~\ref{prop:norm-invariant}, \ref{prop:det};
    Table~\ref{tab:orbits}).
  \item Primitive orbit lengths equal $\pi(m)/2$, half the Pisano period
    (Theorem~\ref{thm:half-pisano}).
  \item The projective ratio $z_t=e_t/(b_t+e_t)$ is attracted to $1/\varphi$ by a
    M\"{o}bius map (Proposition~\ref{prop:projective}).
  \item The harmonic index $\hqa$ has a canonical operator form
    $4\hraw=\mathrm{tr}(\Lambda)$, and the equal-weight Rayleigh identity
    $2\hraw=R(\Lambda,v_0)$ is forced by symmetry
    (Propositions~\ref{prop:trace-rayleigh}, \ref{prop:v0}).
\end{itemize}

These properties are not incidental.  They form a coherent algebraic picture:
the QA system is a mod-$m$ slice of the arithmetic of $\Zfi$, and the harmonic
index $\hqa$ is the natural cross-coupling scalar of this arithmetic.
Applications to convergence certificates and learning-rate optimisation are
developed in the companion paper~\cite{companion}.

%% ────────────────────────────────────────────────────────────────────────────
\bibliographystyle{plain}
\begin{thebibliography}{9}

\bibitem{wall1960}
D.~D. Wall.
\newblock Fibonacci primitive roots.
\newblock \textit{The American Mathematical Monthly}, 67(6):525--532, 1960.

\bibitem{companion}
[Author redacted for review].
\newblock Unified curvature certificates across learning, aggregation,
  attention, arithmetic, and symbolic search.
\newblock \textit{arXiv preprint}, 2026.

\bibitem{hardy2008}
G.~H. Hardy and E.~M. Wright.
\newblock \textit{An Introduction to the Theory of Numbers}, 6th ed.
\newblock Oxford University Press, 2008.

\bibitem{vorobyov1963}
N.~N. Vorobyov.
\newblock \textit{The Fibonacci Numbers}.
\newblock D.~C. Heath, 1963.

\end{thebibliography}

\end{document}
